\documentclass{mimosis}

\usepackage[english]{babel}

\usepackage{graphicx}

% The one colour of links and chapter numbers.
\definecolor{accent}{RGB}{31,58,95}

\usepackage[
    unicode=true,
    plainpages=false,
    pdfpagelabels,
    draft=false,
    hyperfootnotes=false,
    colorlinks=true,
    allcolors={accent},
]{hyperref}

\hypersetup{
    pdfauthor={Author Name},
    pdftitle={Thesis Title},
    pdfsubject={Master's Thesis},
}

% Fonts after classicthesis: Palatino for text, Euler for mathematics, Inconsolata
% for code. One kind of digit everywhere: lining figures, tabular so that columns
% line up, taken by the mathematics from the text so that $16 \to 14$ and \pct
% look like the digits around them.
\usepackage[lining,tabular]{newpxtext}
\usepackage{eulervm}
\usepackage[scaled=0.95,varqu]{zi4}

% Protrusion and font expansion even out the spaces of justified lines; tracking
% spaces the small capitals of headings and running heads.
\usepackage[final,tracking=true]{microtype}
\SetTracking{encoding=*, shape=sc}{50}

% The class sets one and a half line spacing; Palatino reads better tighter.
\setstretch{1.15}

% Headings after classicthesis: the chapter number large and in the colour of the
% links, the running heads in small capitals, captions with a bold label and no
% hanging indent.
\renewcommand*{\chapterformat}{%
  {\fontsize{50}{55}\selectfont\color{accent}\thechapter\autodot\enskip}%
}
\setkomafont{pagehead}{\normalfont\scshape}
\setkomafont{captionlabel}{\usekomafont{caption}\bfseries}
\setcapindent{0pt}

\usepackage[
    backend=bibtex8,
    sortlocale=en_US,
    bibencoding=UTF8,
    maxnames=100,
]{biblatex}
\addbibresource{bibliography.bib}

% Mathematics: symbols, inference rules, and the theorem environments of the
% formal chapters, numbered within each chapter on one counter.
\usepackage{amssymb}
\usepackage{mathtools}
\usepackage{stmaryrd}
\usepackage{mathpartir}
\usepackage{array}

\usepackage[capitalise,nameinlink]{cleveref}

\theoremstyle{plain}
\newtheorem{theorem}{Theorem}[chapter]
\newtheorem{lemma}[theorem]{Lemma}
\newtheorem{proposition}[theorem]{Proposition}
\newtheorem{corollary}[theorem]{Corollary}
\theoremstyle{definition}
\newtheorem{definition}[theorem]{Definition}
\newtheorem{assumption}[theorem]{Assumption}
\newtheorem{example}[theorem]{Example}
\theoremstyle{remark}
\newtheorem{remark}[theorem]{Remark}
\crefname{assumption}{Assumption}{Assumptions}

% Listings of graal.ct, the Cthulhu dialect that models Graal IR.
\usepackage{listings}
\lstdefinelanguage{ct}{
    morekeywords={structure,field,class,constant},
    morecomment=[l]{;},
}
\lstset{
    basicstyle=\ttfamily\footnotesize,
    breaklines=true,
    commentstyle=\itshape,
    columns=fullflexible,
    keepspaces=true,
    upquote=true,
    xleftmargin=1em,
    literate={λ}{{$\lambda$}}1 {→}{{$\to$}}1 {×}{{$\times$}}1 {∅}{{$\emptyset$}}1 {…}{{\ldots}}1,
}
